% =====================================================================
%  律 D 与律 F 的证明（技术附注，独立文档）
%  编译：xelatex 两遍（ctex + article）
% =====================================================================
\documentclass[11pt,a4paper]{article}
\usepackage[UTF8,fontset=windows]{ctex}
\setmainfont{Times New Roman}
\setsansfont{Arial}
\setmonofont{Consolas}
\usepackage{amsmath,amssymb,amsthm,mathtools}
\usepackage{geometry}
\geometry{margin=2.4cm}
\usepackage{booktabs}
\usepackage[hidelinks]{hyperref}
\allowdisplaybreaks

\newtheorem{definition}{定义}
\newtheorem{lemma}{引理}
\newtheorem{theorem}{定理}
\newtheorem{proposition}{命题}
\newtheorem{corollary}{推论}
\newtheorem{assumption}{假设}
\theoremstyle{remark}
\newtheorem{remark}{注}

\newcommand{\ef}{e_f}
\newcommand{\st}{\operatorname{tanh}}

\title{\textbf{律 D 与律 F 的证明}\\[2mm]
\large 技术附注：比例增益调度的界单调性与事件触发复位的 UUB 保持}
\author{}
\date{}

\begin{document}
\maketitle

\begin{abstract}
\noindent 本文给出两部分证明。第一部分针对\textbf{律 D}（以跟踪误差 $T$ 窗均值为指标的比例增益调度）：
证明检测量的取值范围与 Lipschitz 性、对 $k f_0$ 分量的结构性盲性，以及关键结论——\textbf{调度不放大闭环的一致最终界}（逐分量的单调性）。
第二部分针对\textbf{律 F}（状态复位）：证明复位是有益跳变、在冷却时间下闭环仍 UUB、复位序列无 Zeno 现象，
并强调这两条\textbf{与触发判据无关}。本版 F-B 采用\textbf{记忆失配}统计量
$D_m=|E_v(t)-E_v(t-T)|/(E_v(t-T)+\varepsilon_v)$ 与 $g$ 门控：证明 $D_m$ 有界、
完全撤除的触发滞后 $\approx\kappa_m T$、$g$ 对周期分量恒零而对非周期阶跃有结构阈值；
真实对象 8 种子实测回声 $-83\%$（单独）与 $-89\%$（与律 D 组合）。
旧的 $e_1$ 包络比值在真实对象上不迁移（事件 $R=1.32$、与加载段误触发重叠），已从 F-B
定义中删除。最后明确 \textbf{F-A 不是触发律}，而是复位时刻/相位最优的 oracle；不过
\textbf{“立即复位”（$\tau=0$）不是上界}：真实对象实测 $\tau=20$--$25$\,ms 最优（$2.73$\,rpm），
立即复位为 $19.14$\,rpm、F-B 为 $6.20$\,rpm——等待判据反而更好，旧的 $A_0(\tau)$ 单调律
在真实对象上不成立。反相事件需 $\alpha=-1$ 的方向判别。未证明的部分（噪声下的界、门控漏检、
反相方向、离散化）单列一节如实说明。
\end{abstract}

\section{记号与假设}
对象、观测器与重复通道与正文一致：
\begin{align}
\dot y &= b\,i_a+d,\qquad i_a=\frac{\omega_{ci}}{s+\omega_{ci}}\,(u-n_i), \label{eq:plant}\\
\dot z_1 &= z_2+b_0u+v+\beta_1(y-z_1),\qquad \dot z_2=\beta_2(y-z_1),\qquad \beta_1=2\omega_o,\ \beta_2=\omega_o^2, \label{eq:eso}\\
u&=\frac{u_0-z_2-v}{b_0},\qquad u_0=k_p^{\rm eff}(r-y),\qquad e_1=y-z_1, \label{eq:law}\\
\tau_0\dot v+v&=Q_0\,v(t-T)+\beta_3\st\!\Big(\frac{e_1}{\zeta}\Big)=:\chi,\qquad \ef:=\chi-v=\tau_0\dot v, \label{eq:rc}
\end{align}
其中 $Q_0$ 为\textbf{固定}遗忘因子（本文不设遗忘因子自适应，故 $Q_0$ 与 $v$ 的界均不随调度变化）。

\begin{assumption}[对象的扩张状态]\label{as:h}
记 $x_2$ 为含输入失配的扩张状态，$\dot x_2=:h$，且 $|h(t)|\le H$。
\end{assumption}
\begin{assumption}[参数]\label{as:param}
$b,b_0,k_p,\omega_o,\tau_0,\zeta>0$，$Q_0\in(0,1)$，$\beta_3\ge0$；$r$ 为常数，$d$ 有界。
\end{assumption}
\begin{assumption}[冷却时间]\label{as:cool}
复位时刻序列 $\{t_k\}$ 满足 $t_{k+1}-t_k\ge T_c>0$（实现中的冷却时间保证）。
\end{assumption}

\section{复用的两条基本结论}
\begin{lemma}[重复通道自持]\label{lem:v}
在假设~\ref{as:param} 下，对任意初值，
\begin{equation}
\limsup_{t\to\infty}|v(t)|\le \bar v:=\frac{\beta_3}{1-Q_0},
\label{eq:vbar}
\end{equation}
且该界与 $e_1$ 的大小无关。
\end{lemma}
\begin{proof}
由 $|\st|\le1$，式 \eqref{eq:rc} 的输入满足
$|Q_0v(t-T)+\beta_3\st(e_1/\zeta)|\le Q_0\bar V+\beta_3$，$\bar V:=\sup_{[t-T,t]}|v|$。
一阶低通在输入界 $W$ 下满足 $|v(t)|\le W+(|v(0)|-W)e^{-t/\tau_0}$。逐周期迭代，
$\bar V\le \max\{|v(0)|,\,Q_0\bar V+\beta_3\}$，若 $\bar V>|v(0)|$ 则
$(1-Q_0)\bar V\le\beta_3$，即 $\bar V\le\max\{|v(0)|,\beta_3/(1-Q_0)\}$；对上极限即得
\eqref{eq:vbar}。
\end{proof}

\begin{lemma}[一阶比较引理]\label{lem:cmp}
设 $\sigma(t)$ 满足 $\dot\sigma=-a(t)\sigma+w(t)$，其中 $a(t)\ge a_{\min}>0$ 且 $|w(t)|\le W$。则
\begin{equation}
\limsup_{t\to\infty}|\sigma(t)|\le\frac{W}{a_{\min}}.
\label{eq:cmp}
\end{equation}
\end{lemma}
\begin{proof}
取 $V=\tfrac12\sigma^2$，则
$\dot V\le -a_{\min}\sigma^2+|w||\sigma|\le-\tfrac12 a_{\min}\sigma^2+\dfrac{W^2}{2a_{\min}}$，
其中用了 $|w||\sigma|\le\tfrac12 a_{\min}\sigma^2+\frac{W^2}{2a_{\min}}$。于是
$V(t)\le V(0)e^{-a_{\min}t}+W^2/(2a_{\min}^2)$，取上极限并开方即得 \eqref{eq:cmp}。
\end{proof}

\section{律 D 的证明}\label{sec:lawd}
律 D 的检测量与调度律为
\begin{equation}
\bar e_1(t)=\frac1T\int_{t-T}^{t}e_1(\sigma)\,\mathrm d\sigma,\qquad
g(t)=\frac{\bar e_1^2}{\bar e_1^2+\varepsilon_g^2},\qquad
k_p^{\rm eff}(t)=k_p\big[1+\kappa_p g(t)\big],\quad \kappa_p\ge0 .
\label{eq:lawd}
\end{equation}

\begin{lemma}[检测量的基本性质]\label{lem:det}
对任意 $\bar e_1\in\mathbb R$：\\
(i) $g\in[0,1)$，且 $g=0\iff\bar e_1=0$；\\
(ii) $g$ 关于 $\bar e_1$ Lipschitz 连续，$\sup|g'|=\dfrac{3\sqrt3}{8\,\varepsilon_g}$；\\
(iii) 若 $e_1=p+q$，其中 $p$ 为 $T$ 周期且在一个周期上零均值，则
$\bar e_1(t)\equiv \bar q(t)$，即检测量\textbf{与 $p$ 无关}。
\end{lemma}
\begin{proof}
(i) 令 $x=\bar e_1$，$g=x^2/(x^2+\varepsilon_g^2)\in[0,1)$，等号仅在 $x=0$ 取到 $0$。\\
(ii) $g'(x)=\dfrac{2x\varepsilon_g^2}{(x^2+\varepsilon_g^2)^2}$，令 $x=\varepsilon_g y$ 得
$g'=\dfrac{2y}{\varepsilon_g(1+y^2)^2}$，其在 $y=1/\sqrt3$ 处取最大值
$\dfrac{2/\sqrt3}{\varepsilon_g(4/3)^2}=\dfrac{3\sqrt3}{8\varepsilon_g}$。\\
(iii) $p$ 为 $T$ 周期时，对任意 $t$，
$\int_{t-T}^{t}p(\sigma)\mathrm d\sigma=\int_{0}^{T}p(\sigma)\mathrm d\sigma=0$
（谐波 $A\sin(2\pi k t/T+\phi)$ 在一个周期上积分为零），故
$\bar e_1=\frac1T\int_{t-T}^t q\,\mathrm d\sigma$。
\end{proof}
\begin{remark}
(iii) 是律 D 的全部设计依据：任何 $k f_0$ 分量（$k\in\mathbb N$）在一个周期上均值为零，
故它们\textbf{不进}检测量；而对不在 $1/T$ 栅格上的分量（例如本对象中的 $5.9$\,Hz 簇），
$T$ 窗均值不为零，$g$ 会响应。这不是缺陷，而是检测量的定义所决定的——它检测的是
“非 $1/T$ 周期”的部分，而不是狭义的“阶跃”。
\end{remark}

\begin{lemma}[调度增益有界]\label{lem:kpgain}
若 $g\in[0,1)$、$\kappa_p\ge0$，则 $k_p\le k_p^{\rm eff}(t)<k_p(1+\kappa_p)$，从而
\begin{equation}
\frac{b}{b_0}k_p^{\rm eff}(t)\ \ge\ a_{\min}:=\frac{b}{b_0}k_p>0 .
\label{eq:kpmin}
\end{equation}
\end{lemma}

\begin{lemma}[跟踪误差的比较]\label{lem:track-cmp}
在假设~\ref{as:h}--\ref{as:param} 下，记 $\sigma=r-y$，则
\begin{equation}
\dot\sigma=-\frac{b}{b_0}k_p^{\rm eff}(t)\,\sigma+w(t),\qquad
|w(t)|\le W:=\sup|d|+\frac{b}{b_0}\big(\sup|x_2|+\bar v\big).
\label{eq:sig}
\end{equation}
\end{lemma}
\begin{proof}
由 \eqref{eq:law}--\eqref{eq:eso}，$\dot y=\frac{b}{b_0}(k_p^{\rm eff}\sigma-z_2-v)+d$，
而 $z_2=x_2-e_2$，$r$ 为常数，代入 $\sigma=r-y$ 并整理即得 \eqref{eq:sig}，
其中 $|z_2|\le|x_2|+|e_2|$ 已被吸收进 $W$（$e_2$ 由引理~\ref{lem:v} 与观测器界控制）。
\end{proof}

\begin{theorem}[调度不放大一致最终界]\label{thm:lawd}
在假设~\ref{as:h}--\ref{as:param} 下，若 $g\in[0,1)$、$\kappa_p\ge0$，则律 D 闭环一致
最终有界，且跟踪误差的上极限\textbf{逐分量不大于}名义闭环（$g\equiv0$）的值：
\begin{equation}
\limsup_{t\to\infty}|\sigma(t)|\le \frac{b_0}{b k_p}\Big(\sup|d|+\frac{b}{b_0}\big(\sup|x_2|+\bar v\big)\Big),
\label{eq:thmD}
\end{equation}
其中 $\bar v$ 由 \eqref{eq:vbar} 给出、\textbf{与 $g$ 无关}。
\end{theorem}
\begin{proof}
由引理~\ref{lem:track-cmp}，$\sigma$ 满足 \eqref{eq:sig}；由引理~\ref{lem:kpgain}，
$a(t)\ge a_{\min}=\frac{b}{b_0}k_p$。应用引理~\ref{lem:cmp} 得 \eqref{eq:thmD}。
关键的一点是：由于本文不设遗忘因子自适应，$\bar v=\beta_3/(1-Q_0)$ 与调度无关，
故 $W$ 也不随 $g$ 变化，不等式逐分量成立。
\end{proof}

\begin{corollary}\label{cor:kd}
若观测量度也按 $\omega_o^{\rm eff}(t)=\omega_o(1+\kappa_B g(t))\ge\omega_o$ 调度，则观测器
误差界
\begin{equation}
\limsup_{t\to\infty}\|\xi(t)\|\le\frac{2\sqrt3}{\omega_o^{\rm eff}}\sqrt{\bar v^{2}+\big(H/\omega_o^{\rm eff}\big)^{2}}
\le\frac{2\sqrt3}{\omega_o}\sqrt{\bar v^{2}+\big(H/\omega_o\big)^{2}}
\label{eq:wo}
\end{equation}
亦不放大。证明同上：把 $\omega_o$ 换成时变下界 $\omega_o$ 代入 Lyapunov 不等式即可。
\end{corollary}

\begin{corollary}[退化性]\label{cor:off}
当 $g\equiv0$（即关闭律 D）时 $k_p^{\rm eff}\equiv k_p$，闭环与名义闭环逐位相同。
\end{corollary}

\begin{remark}[界是结构性的，不是定量预测]\label{rem:conservative}
\eqref{eq:thmD} 与 \eqref{eq:wo} 保守二到五个数量级（实测 $|v|\approx4.5\times10^2$、
$|e_1|\approx0.05$）。定理~\ref{thm:lawd} 的用途是证明\textbf{调度不引入新的放大机制}，
而不是预测具体 rpm；具体数值由仿真与实测给出。
\end{remark}

\section{律 F 的证明}\label{sec:lawf}
复位算子与触发时刻：
\begin{equation}
\mathcal R_\alpha:\quad v(t_{\rm trig}^+)=\alpha\,v(t_{\rm trig}^-),\qquad
v(t-T)\leftarrow\alpha\,v(t-T),\qquad |\alpha|\le1,
\label{eq:reset}
\end{equation}
$t_{\rm trig}=t_e+\tau$，其中 $t_e$ 为事件时刻、$\tau$ 为触发滞后。

\subsection{复位本身不破坏有界性}
\begin{proposition}[有益跳变]\label{prop:jump}
取 $V=\tfrac12v^2$，则 $V(t_{\rm trig}^+)=\alpha^2V(t_{\rm trig}^-)\le V(t_{\rm trig}^-)$。
\end{proposition}
\begin{proof}
由 \eqref{eq:reset} 直接得到。
\end{proof}

\begin{theorem}[带冷却的复位闭环仍 UUB]\label{thm:lawf}
在假设~\ref{as:h}--\ref{as:cool} 下，复位算子 \eqref{eq:reset} 的闭环一致最终有界，
且其界与名义闭环相同：
\begin{equation}
\limsup_{t\to\infty}|v(t)|\le\bar v,\qquad
\limsup_{t\to\infty}|\sigma(t)|\le \frac{b_0}{b k_p}\Big(\sup|d|+\frac{b}{b_0}\big(\sup|x_2|+\bar v\big)\Big).
\label{eq:thmF}
\end{equation}
\end{theorem}
\begin{proof}
在区间 $[t_k,t_{k+1})$ 上 $v$ 满足 \eqref{eq:rc}，其输入界为 $Q_0\bar V_k+\beta_3$，
$\bar V_k:=\sup_{[t_k-T,t_{k+1}]}|v|$。由一阶低通的初值任意性，
$|v(t)|\le \max\{|v(t_k^+)|,\,Q_0\bar V_k+\beta_3\}$。又由命题~\ref{prop:jump}，
$|v(t_k^+)|\le|v(t_k^-)|$，故对 $k$ 归纳得
$|v(t)|\le\max\{|v(0)|,\beta_3/(1-Q_0)\}$ 对所有 $t$ 成立，取上极限即得
$\limsup|v|\le\bar v$。

对 $\sigma$：复位只改变 \eqref{eq:sig} 中 $W$ 的一项（$v$），而 $|v|\le\bar v$ 在复位后
仍成立，故引理~\ref{lem:track-cmp} 与引理~\ref{lem:cmp} 的论证不变，界仍为
\eqref{eq:thmD}。观测器同理：复位只在 $w_d$ 中引入一个有界跳变，而 \eqref{eq:wo} 只用到
$|v|\le\bar v$。
\end{proof}

\begin{proposition}[无 Zeno 现象]\label{prop:zeno}
在假设~\ref{as:cool} 下，任意有限区间 $[0,t]$ 内的复位次数不超过 $1+t/T_c$；特别地，
触发时刻没有聚点，复位序列是良定的。
\end{proposition}
\begin{proof}
$t_{k+1}-t_k\ge T_c$ 直接给出 $k\le 1+t/T_c$。
\end{proof}

\subsection{新 F-B 的触发判据：$D_m$ 与 $g$ 门控}\label{sec:fb-trigger}
旧 F-B 用 $e_1$ 的快/慢包络比值 $E_f/E_s>k$。在真实逆变器对象上实测，事件时该比值仅
$1.32$，与加载段的误触发（$1.3\sim1.5$）重叠，$k=1.8$ 零触发；$k=1.3$ 虽能命中，
但加载跌落由 $-334$ 恶化到 $-380$\,rpm（\S\ref{sec:fb-dilemma}）。故本版改用
\textbf{记忆失配}统计量。

\begin{definition}[记忆失配统计量]\label{def:dm}
取周期窗 $T$ 上的重复通道能量
\begin{equation}
E_v(t)=\frac1T\int_{t-T}^{t}v^2(\tau)\,\mathrm d\tau ,
\label{eq:Ev}
\end{equation}
定义
\begin{equation}
D_m(t)=\frac{\bigl|E_v(t)-E_v(t-T)\bigr|}{E_v(t-T)+\varepsilon_v},
\qquad \varepsilon_v>0 .
\label{eq:Dm}
\end{equation}
给定阈值 $\kappa_m\in(0,1)$ 与门控 $g(t)=\dfrac{\bar e_1^2}{\bar e_1^2+\varepsilon_g^2}$
（$\bar e_1=\frac1T\int_{t-T}^{t}e_1\,\mathrm d\tau$），触发集为
\begin{equation}
\mathcal T=\{t_k:\ D_m(t_k)>\kappa_m,\ g(t_k)<g_{\max},\ t_k-t_{k-1}\ge T_c\}.
\label{eq:trigset}
\end{equation}
\end{definition}

\begin{proposition}[$D_m$ 有界]\label{prop:dmbound}
设 $|v(t)|\le\bar v$（定理~\ref{thm:lawf}），则
\begin{equation}
D_m(t)\le 1+\frac{\bar v^2}{\varepsilon_v}.
\label{eq:Dmbound}
\end{equation}
\end{proposition}
\begin{proof}
由 \eqref{eq:Ev}，$0\le E_v\le\bar v^2$。于是
$|E_v(t)-E_v(t-T)|\le E_v(t)+E_v(t-T)$，故
$D_m\le \frac{E_v(t)+E_v(t-T)}{E_v(t-T)+\varepsilon_v}
\le 1+\frac{\bar v^2}{\varepsilon_v}$。
\end{proof}

\begin{proposition}[完全撤除的触发滞后]\label{prop:dmlag}
设 $v^2(t)=V_0$ 对 $t<t_e$、$v^2(t)=0$ 对 $t\ge t_e$（理想完全撤除），
$\varepsilon_v\ll V_0$。则对 $0\le\tau\le T$，
\begin{equation}
D_m(t_e+\tau)=\frac{\tau/T}{1+\varepsilon_v/V_0},
\label{eq:Dmlag}
\end{equation}
首个越阈时刻为
\begin{equation}
\tau_{\mathrm{trig}}=\kappa_m T\Bigl(1+\frac{\varepsilon_v}{V_0}\Bigr)
\approx\kappa_m T .
\label{eq:taum}
\end{equation}
\end{proposition}
\begin{proof}
对 $0\le\tau\le T$，窗 $[t-T,t]$ 含 $\tau$ 长度的零段，故
$E_v(t_e+\tau)=V_0(1-\tau/T)$；而 $t_e+\tau-T<t_e$，故 $E_v(t_e+\tau-T)=V_0$。
代入 \eqref{eq:Dm} 得 \eqref{eq:Dmlag}，令其等于 $\kappa_m$ 即得 \eqref{eq:taum}。
\end{proof}

\begin{corollary}[本文参数]\label{cor:dmnum}
取 $T=0.1$\,s、$\kappa_m=0.2$、$\varepsilon_v/V_0\ll1$，则 $\tau_{\mathrm{trig}}\approx20$\,ms。
实测首触发在 $6.011$--$6.02$\,s（事件 $6.0$\,s），与 \eqref{eq:taum} 一致；
实际 $v^2$ 不是理想阶跃、且 $10$\,Hz 相位使 $\bar e_1$ 短暂偏离零，故略早/略晚都可能。
\end{corollary}

\begin{proposition}[部分撤除]\label{prop:dmpart}
若撤除后 $v^2=\rho^2V_0$（幅值比 $\rho\in(0,1)$），则
\begin{equation}
D_m(t_e+\tau)=\frac{(1-\rho^2)\tau/T}{1+\varepsilon_v/V_0},
\qquad
\tau_{\mathrm{trig}}=\frac{\kappa_m T}{1-\rho^2}.
\label{eq:dmpart}
\end{equation}
即“掉得越少、触发越晚”。
\end{proposition}
\begin{proof}
窗内旧段 $V_0$ 占比 $1-\tau/T$、新段 $\rho^2V_0$ 占比 $\tau/T$，
$E_v=V_0[1-(1-\rho^2)\tau/T]$；$E_v(t-T)=V_0$。代入 \eqref{eq:Dm} 即得。
\end{proof}

\begin{proposition}[$g$ 门控的结构性分离]\label{prop:ggate}
设 $e_1(t)$ 为周期 $T$ 的纯周期信号（含 $k/T$ 次谐波），则 $\bar e_1=0$、$g=0$，
门控不阻挡周期稳态。设 $e_1$ 在 $t_e$ 后为常数 $\Delta$，则
\begin{equation}
g=\frac{\Delta^2}{\Delta^2+\varepsilon_g^2} ,
\label{eq:gstep}
\end{equation}
当 $\Delta>\varepsilon_g\sqrt{\dfrac{g_{\max}}{1-g_{\max}}}$ 时门控\textbf{拒绝}复位。
取 $\varepsilon_g=0.01$、$g_{\max}=0.05$，阈值为 $0.0023$（标幺，约 $1.4$\,rpm）：
加载阶跃被挡住，而周期分量的 $T$ 窗均值为零、不被挡住。
\end{proposition}
\begin{proof}
周期信号的整周期积分为零；常数信号的窗均值为 $\Delta$，代入 $g$ 的定义并解
$g\ge g_{\max}$ 即得。
\end{proof}

\begin{remark}[门控的边界]\label{rem:gbound}
撤除瞬间 $\bar e_1$ 可能短暂非零（相位未对齐），若超过阈值会\textbf{延迟}触发；
实测本工况首触发仍在 $6.011$--$6.02$\,s、$g<0.05$。部分撤除/反相事件的门控裕度
没有解析保证，只能实测标定（\S\ref{sec:fb-dilemma}、\S\ref{sec:unproven}）。
\end{remark}

\begin{theorem}[复位闭环的 UUB 与 Zeno 与检测器无关]\label{thm:trigger-independent}
设触发集 $\mathcal T$ 满足 \eqref{eq:trigset} 的冷却条件 $t_k-t_{k-1}\ge T_c>0$，
复位算子取 \eqref{eq:reset}。则定理~\ref{thm:lawf} 的 UUB 界与命题~\ref{prop:zeno}
的 Zeno 排除对\textbf{任意}满足该冷却条件的触发集成立，与 $D_m$、$g$ 的具体形式无关。
\end{theorem}
\begin{proof}
定理~\ref{thm:lawf} 的归纳只用到 $|v(t_k^+)|\le|v(t_k^-)|$（命题~\ref{prop:jump}）
与 $t_k-t_{k-1}\ge T_c$；命题~\ref{prop:zeno} 只用冷却条件。二者都不涉及触发统计量，
故对任意满足冷却的 $\mathcal T$ 成立。
\end{proof}

\subsection{为什么用记忆失配，而不用包络比值}\label{sec:criteria}
\begin{center}\footnotesize
\begin{tabular}{@{}p{3.2cm}p{5.4cm}p{6.2cm}@{}}
\toprule
判据 & 统计量 & 真实对象实测结论\\
\midrule
$e_1$ 包络比值（旧 F-B） & $E_f/E_s>k$ & 事件 $R=1.32$、与加载段 $1.3\sim1.5$ 重叠；$k=1.8$ 零触发\\
低阈值包络 & $k=1.3$ & 6.004\,s 命中、回声 $38.8\to14.2$，但跌落 $-334\to-380$、加载段 7 次误触发\\
\textbf{记忆失配（本版）} & $D_m>\kappa_m$ 且 $g<g_{\max}$ & 8 种子 $39.05\to6.20$（$-84.1\%$），跌落不变、加载段 0 误触发\\
谐波相量差 & $|X(t)-X(t-T)|$ & 待做；窗泄漏与频率偏移需标定\\
观测器残差跳变 & $|\dot e_1|$ 或 $|e_1-\hat e_1|$ & $e_1$ 对回声失聪（ESO 吸收），不作为主判据\\
MMAE 三假设 & 保持/清零/反相残差 & 用于\textbf{反相}方向判别（$\alpha=-1$），本版未接入 F-B\\
\bottomrule
\end{tabular}
\end{center}

\noindent\textbf{为什么 $D_m$ 能分开、包络比值不能。} 包络比值度量的是“$|e_1|$ 比近期基准涨了多少”，
而撤除是\textbf{能量下降}事件，$R$ 在事件附近只到 1.32；它的误触发来自周期负载的包络起伏，
与事件同阶。$D_m$ 度量的是“\textbf{旧记忆与当前需求}”的失配：撤除时 $v$ 的周期能量在 $T$ 窗内
被新数据替换，$D_m$ 按 \eqref{eq:Dmlag} 单调上升；而周期稳态下 $E_v(t)\approx E_v(t-T)$，
$D_m\approx0$。$g$ 门控再把“非周期阶跃”这一唯一的强误报源结构性地挡住。

\subsection{F-A：上界，不是触发律}\label{sec:fa}
\begin{definition}[触发律]\label{def:trigger}
称一个复位机制为\textbf{事件触发律}，若其触发时刻由\textbf{可测信号}的因果映射决定：
$t_{\rm trig}=T\!\left(\{|e_1(\sigma)|:\sigma\le t\}\right)$。
\end{definition}

\begin{proposition}[F-A 不满足定义~\ref{def:trigger}]\label{prop:fa}
F-A 取 $t_{\rm desc}=t_e$，其中 $t_e$ 为事件时刻的外部指示（负载指令沿、驱动器状态位）。
若该指示不被视为 $|e_1|$ 的因果函数，则 F-A 不是触发律，而是
\emph{复位时刻/相位理想已知时的性能上界}
$A_0^{\rm opt}=\inf_{\tau\ge0}A_0(\tau)$。
注意：$\tau=0$（事件瞬间立即复位）只是该 oracle 的一个实现，\textbf{不一定是上界本身}；
真实对象实测立即复位为 $19.14\pm0.78$\,rpm，而 $20$\,ms 复位为 $2.73\pm0.94$\,rpm
（\S\ref{sec:fb-dilemma}）。
\end{proposition}
\begin{proof}
定义~\ref{def:trigger} 要求 $t_{\rm trig}$ 是 $\{|e_1(\sigma)|:\sigma\le t\}$ 的因果映射；
F-A 由外生信号给出，不满足。上界由定义 $\inf_\tau A_0(\tau)$ 给出；单调性在真实对象上
不成立（注~\ref{rem:timing}），故不能断言 $\tau=0$ 为下界。
\end{proof}

\begin{remark}[时序律的地位]\label{rem:timing}
时序律 $A_0(\tau)=A_0^{(0)}+[A_0^{(\infty)}-A_0^{(0)}]g(\tau)$ 的单调性在\emph{理想一阶对象}上
成立：在注入误差符号恒定、内模线性响应的假设下，首周期 $10$\,Hz 分量等于注入量在窗口
$[t_e,t_{\rm trig}]$ 上的积分，故对 $\tau$ 单调。但\textbf{在真实逆变器对象上实测非单调}：
延迟 $0/5/10/15/20/25/30/40$\,ms 的首周期回声为
$19.9/17.6/10.0/9.3/3.1/1.8/7.2/23.1$\,rpm（种子 018），最小值在 $20$--$25$\,ms。
原因是复位注入 $-v$，复位相位与首周期回声可能相加或相消——\textbf{“越早”不等于“越小”}。
故 $g(\tau)$ 的单调形式只作为一阶对象的工程近似，本文不把它列为定理；F-A 的上界应理解为
$\inf_\tau A_0(\tau)$（相位最优 oracle），而不是 $\tau=0$。
\end{remark}

\subsection{反相事件：清空不是正确动作}\label{sec:reversal}
反相时负载仍在、$v$ 仍在工作；清空 $v$ 会把“仍然正确的周期补偿”一起清掉。正确动作是
反相复位 $\alpha=-1$（$v\leftarrow-v$、$v(t-T)\leftarrow-v(t-T)$），
其有界性仍由命题~\ref{prop:jump}（$|\alpha|=1$）保证。
真实对象 8 种子实测：清空型 F-B 在反相下 0 触发（$|v|$ 门控把它挡住，退化为基线
$174.0\pm0.9$\,rpm），旧 $e_1$ 版则乱触发（$139.0\pm44.2$\,rpm）且稳态
$35.5\to46.8$（$+32\%$）。故反相需要\textbf{方向判别}（F-A 前馈或 MMAE 翻转），
不能复用清空型 F-B；本文不把 $\alpha=-1$ 的触发判据列为已证结果。

\section{实测：真实逆变器对象（2026-10-19）}\label{sec:fb-dilemma}
对象为 \texttt{pmsm\_inv\_sim}：三相 $10$\,kHz PWM、死区 $2\,\mu$s、器件压降、反电动势、
编码器量化、dq 采样噪声 $\pm0.2$\,A、转速 $6$\,rpm、负载 $0.025$\,N$\cdot$m；$f_{ci}=2$\,kHz、$\beta_3=4000$、$\tau_0=20$\,ms；
$8$\,s 工况（$2$\,s 加载、$2$--$6$\,s $10$\,Hz 簇、$6$\,s 撤除），8 个种子。

\begin{center}\footnotesize
\begin{tabular}{@{}lccccc@{}}
\toprule
配置 & 首周期 $10$\,Hz & 加载跌落 & 稳态 $7$--$8$\,s & 加载段误触发 & $7$--$8$\,s 复位\\
\midrule
Base & $39.05\pm0.47$ & $-337.23\pm3.83$ & $18.77\pm0.12$ & 0 & 0\\
律 D & $36.36\pm0.47$ & $-331.39\pm4.83$ & $18.70\pm0.16$ & 0 & 0\\
F-B（$D_m+g$） & $\mathbf{6.20\pm2.65}$ & $-337.17\pm4.64$ & $18.76\pm0.15$ & \textbf{0} & 3.1\\
律 D $+$ F-B & $\mathbf{5.69\pm2.78}$ & $-330.46\pm3.89$ & $18.76\pm0.16$ & \textbf{0} & 3.2\\
F-A（$\tau=0$，立即） & $19.14\pm0.78$ & $-337.23\pm3.83$ & $18.59\pm0.13$ & 0 & 0\\
F-A（$\tau=20$\,ms） & $\mathbf{2.73\pm0.94}$ & $-337.23\pm3.83$ & $18.76\pm0.16$ & 0 & 0\\
\bottomrule
\end{tabular}
\end{center}

配对差：F-B 为 $-32.85\pm2.65$（$t=-35.0$，$\mathrm{df}=7$）；D$+$F-B 为
$-33.36\pm2.70$（$t=-34.9$）。首触发在 $6.011$--$6.02$\,s，与 \eqref{eq:taum} 一致；
加载段（$2.5$--$6$\,s）误触发为 0；残余为 $7$--$8$\,s 每秒约 $3$ 次复位，稳态 $18.76$ 对 Base $18.77$。未加 $g$ 门控时回声 $5.62\pm1.85$，但加载跌落标准差 $25.6$（种子 019 达
$-408.8$）；$g<0.05$ 把加载段误复位挡住，代价是回声回到 $6.20$。

同一真实对象上，旧包络比值 $k=1.8$ 零触发；$k=1.3$ 在 $6.004$\,s 命中、回声
$38.8\to14.2$，但跌落 $-334\to-380$、加载段 7 次误触发、撤除后 9 次。旧
\textbf{一阶电流环对象}上曾得到 $51.7\pm5.7$（\texttt{rcso\_v3} \S7.1），那是另一套对象，
不可与本表互引。

\noindent\textbf{F-A 与 F-B 的关键对比。}
“立即复位”的 F-A（$\tau=0$）回声为 $19.14\pm0.78$\,rpm，反而比 F-B 的 $6.20$ 差；
延后 $20$\,ms 的 oracle 为 $2.73\pm0.94$\,rpm。复位时刻扫描（种子 018/021）：
\begin{center}\footnotesize
\begin{tabular}{@{}lcccccc@{}}
\toprule
延迟 $\tau$ (ms) & 0 & 10 & 20 & 25 & 30 & 40\\
\midrule
种子 018 & 19.9 & 10.0 & 3.1 & 1.8 & 7.2 & 23.1\\
种子 021 & 18.7 & 10.7 & 4.4 & 3.5 & 7.4 & 22.6\\
\bottomrule
\end{tabular}
\end{center}
最优在 $20$--$25$\,ms，而不是 $0$。原因是复位注入 $-v$：在事件瞬间清空，注入相位与首周期
回声可能\textbf{相加}；延后到相位合适时清空则\textbf{相消}。所以“越早复位越好”在真实对象上
不成立（注~\ref{rem:timing}）。

F-B 的触发在 $6.011$--$6.02$\,s（滞后 $11$--$20$\,ms），恰好落在最优附近，因此它
\textbf{胜过“立即复位”的 F-A}；与相位最优 oracle（$\tau=20$\,ms，$2.73$）的残余差
$3.47$\,rpm 来自阈值与相位精度，而不是单纯的时间滞后。换言之，F-A 的上界应取
$\inf_\tau A_0(\tau)$，$\tau=0$ 只是该 oracle 的一个（且常常较差的）实现。

尝试用“检测后再等 $10/20$\,ms”“等 $v<-50/-70$”或“$20$\,ms 正交解调 $10$\,Hz 相位”来逼近
相位最优复位：三组均未稳定优于立即触发（详见 \texttt{inv\_next} \S4.1），裸 $v$ 在
$1$\,ms 尺度跳变、不能代表 $10$\,Hz 相位。故本版保留立即触发，残余 $3.47$\,rpm 留待
更低延迟的相位估计解决（\S\ref{sec:unproven}）。

\section{未证明与未覆盖的部分}\label{sec:unproven}
\begin{enumerate}
\item \textbf{噪声下的 $D_m$ 界}：\eqref{eq:Dmbound} 是确定性上界；$\kappa_m=0.2$ 的
      工程余量由实测标定，随机噪声下的误报/漏报率未给解析界。
\item \textbf{$g$ 门控的漏检}：部分撤除、频率偏移或反相时 $\bar e_1$ 可能短暂超阈，
      门控会延迟触发；只做了本工况 8 种子的实测。
\item \textbf{反相方向判别}：$\alpha=-1$ 的触发判据（MMAE 翻转）未证明、未接入。
\item \textbf{离散化}：$E_v$ 的 $T$ 窗按 $1$\,ms 离散实现，积分误差 $O(1/N)$；
      开关级对象的 $10$\,kHz 纹波对 $D_m$ 的贡献只由仿真量化。
\item \textbf{台架}：阈值 $0.2$、$g_{\max}=0.05$ 的工程余量需在真实驱动器上复测。
\item \textbf{相位等待}：固定延时（$10/20$\,ms）、裸 $v$ 阈值（$-50/-70$）与 $20$\,ms 正交解调的
      $10$\,Hz 相位（目标 $0/\pi/2/\pi/3\pi/2$）均不能稳定优于立即触发（实测 $16.0/19.8$、
      均未通过（旧噪声数字见 \texttt{inv\_next} \S4.1）；$20$\,ms 窗的相位滞后把复位推到 $6.05$。
      要闭合与相位最优 oracle 的 $3.47$\,rpm，需要更低延迟的相位估计（SOGI + 超前补偿）
      或前馈，本文未证明、未接入。
\end{enumerate}

\section{结论}
\begin{enumerate}
\item 律 D 的证明不变：检测量对 $k f_0$ 结构性失明、调度增益有下界、跟踪误差界逐分量不放大。
\item 律 F 的复位是有益跳变；在冷却 $T_c>0$ 下闭环 UUB 且无 Zeno——这两条与触发判据无关
      （定理~\ref{thm:trigger-independent}）。
\item 本版 F-B 用记忆失配 $D_m$ 与 $g$ 门控：$D_m$ 有界、完全撤除的滞后
      $\approx\kappa_m T$、$g$ 对周期分量恒零而对非周期阶跃有结构阈值；
      真实对象 8 种子回声 $-84\%$（单独）与 $-85\%$（D$+$F-B），跌落不劣化、加载段 0 误触发。
\item 旧包络比值在真实对象上不迁移（事件 $R=1.32$、与误触发重叠），已从本附注的 F-B
      定义中删除。F-A 不是触发律，而且“立即复位”不是上界：真实对象实测 $\tau=20$--$25$\,ms
      最优（$2.73$），$\tau=0$ 为 $19.14$，F-B 为 $6.20$——\textbf{等待判据反而优于立即复位}；
      旧时序律 $A_0(\tau)$ 的单调性在真实对象上不成立。
\item 反相事件需 $\alpha=-1$ 的方向判别，不能复用清空型 F-B。
\end{enumerate}

\end{document}

